\documentclass[../main.tex]{subfiles}
\begin{document}

\section{Lecture 16 -- Potential Flow Around a Sphere and D'Alembert's Paradox}\label{sec:lec16}
\begin{center}
  \textbf{Date:} June 1, 2026
\end{center}

\subsection*{Overview}
Lecture~15 reduced ideal, incompressible, irrotational flow to a single linear field equation, $\lapl\phi=0$ (\cref{sec:lec15:incompressible}). This lecture exhibits the single most important closed-form solution of that program: the flow generated by a \emph{sphere moving through a fluid otherwise at rest}. The payoff is threefold. First, the boundary-value problem turns out to be mathematically identical to the electrostatics of a conducting sphere in a uniform field, so the solution is a \emph{pure dipole}. Second, integrating the (unsteady) Bernoulli pressure over the sphere yields a force that is \emph{exactly zero whenever the sphere moves at constant velocity} -- this is \textbf{D'Alembert's paradox}. Third, when the sphere accelerates, the surviving term reveals that the fluid endows the sphere with an \textbf{added (effective) mass} equal to half the displaced fluid mass. We close by promoting the effective mass to a symmetric rank-2 tensor for general bodies, and by previewing how viscosity resolves the paradox.

\medskip\noindent\textbf{Topics:}
\begin{enumerate}
  \item Recap of the potential-flow program (Laplace $+$ Bernoulli-like pressure)
  \item Physical setup: a sphere moving through a fluid at rest; two reference frames
  \item Qualitative expectations: streamlines, stagnation points, the drag question
  \item Boundary-value problem in the sphere frame: Laplace $+$ Neumann at $r=R$ $+$ dipole at infinity
  \item Multipole expansion; why the velocity BC selects the pure $\ell=1$ dipole
  \item Matching the Neumann condition at $r=R$; solving for the dipole strength
  \item Galilean transformation to the fluid frame
  \item Streamline / equipotential analysis and the stagnation points
  \item Generalization to arbitrary motion $\bvec{r}_s(t)$
  \item The force from the unsteady Bernoulli pressure; parity and the vanishing of drag
  \item Added mass: factor $\tfrac{1}{2}$ for the sphere; the effective-mass tensor $M_{ij}$
  \item D'Alembert's paradox and its resolution by viscosity
\end{enumerate}

%%─────────────────────────────────────────────────────────────────────────────
\subsection{Recap: The Potential-Flow Program}\label{sec:lec16:recap}
%%─────────────────────────────────────────────────────────────────────────────

\subsubsection*{Conceptual Background}
Before tackling the sphere, recall precisely \emph{why} the flow problem became tractable. We made a chain of assumptions, each linearising or decoupling the dynamics one step further.

\paragraph{Step 1 -- Irrotational $\Rightarrow$ scalar potential.}
We defined the \emph{vorticity} $\bvec{\omega}\equiv\curl\vel$. For an ideal (inviscid) fluid, Kelvin's circulation theorem (Lec~14) guarantees that if the flow is irrotational far from any obstacle -- as it is for a body moving through a fluid at rest -- it remains irrotational everywhere the fluid visits. A curl-free field is a gradient, so
\begin{equation}
  \bvec{\omega}=\curl\vel=\bvec{0}\;\Longrightarrow\;\vel=\grad\phi,
  \label{eq:lec16_vel_pot}
\end{equation}
the identity $\curl\grad\phi\equiv\bvec{0}$ following from the commutativity of partial derivatives $\epsilon_{ijk}\partial_j\partial_k\phi=0$. A three-component vector field is replaced by one scalar.

\paragraph{Step 2 -- Incompressible $\Rightarrow$ Laplace.}
For a constant-density fluid the continuity equation collapses to $\divg\vel=0$. Substituting \eqref{eq:lec16_vel_pot},
\begin{equation}
  \divg\vel=\divg(\grad\phi)=\lapl\phi=0.
  \label{eq:lec16_laplace}
\end{equation}
\textit{Significance:} The full \emph{nonlinear} Euler problem has collapsed onto the \emph{linear} Laplace equation. Linearity is the crucial gain: it licenses superposition, which is exactly what makes the multipole method below work.

\paragraph{Step 3 -- Pressure from the unsteady Bernoulli relation.}
Pressure is no longer a dynamical unknown; it is read off after $\phi$ is known. The Bernoulli-like equation for incompressible irrotational flow (Lec~15) is
\begin{equation}
  \boxed{\pdv{\phi}{t} + \tfrac{1}{2}v^2 + \frac{p}{\rho} + \Phi_g = f(t),}
  \label{eq:lec16_bernoulli}
\end{equation}
where $\Phi_g$ is the gravitational potential and $f(t)$ an arbitrary function of time. We will drop gravity ($\Phi_g=0$) throughout, since for aerodynamic/hydrodynamic drag it is irrelevant (or, for ships, treated separately as buoyancy). Solving for the pressure:
\begin{equation}
  \boxed{p = \rho f(t) - \rho\,\pdv{\phi}{t} - \tfrac{1}{2}\rho v^2.}
  \label{eq:lec16_pressure}
\end{equation}
\textit{Significance:} Each term has a clean meaning. $\rho f(t)$ is a spatially uniform reference pressure set by the far field; $-\rho\,\partial_t\phi$ is the genuinely \emph{unsteady} contribution, present only when the flow changes in time; $-\tfrac{1}{2}\rho v^2$ is the familiar Bernoulli suction -- fast fluid is low-pressure fluid. The whole force calculation later hinges on which of these terms survive integration over the sphere.

\nt{The program is asymmetric: \eqref{eq:lec16_laplace} is \emph{one equation in one unknown} $\phi$ and contains no time derivative -- time enters \emph{only} through the boundary conditions. Once $\phi(\curcoord,t)$ is found, \eqref{eq:lec16_pressure} delivers the pressure everywhere. This ``solve-then-read-off'' structure is what makes the moving-sphere problem soluble in closed form.}

%%─────────────────────────────────────────────────────────────────────────────
\subsection{Physical Setup and Two Reference Frames}\label{sec:lec16:setup}
%%─────────────────────────────────────────────────────────────────────────────

\subsubsection*{Physical Motivation}
Consider a rigid sphere of radius $R$ moving with velocity $\bvec{v}_0$ through an unbounded incompressible fluid that is at rest at infinity. This is the prototype of every engineering problem of motion through a fluid: a submarine or ship hull through water, an aircraft or projectile through air. We choose a sphere not because real vehicles are spherical -- they are streamlined precisely to \emph{minimise} drag and avoid turbulent wakes -- but because the sphere is the simplest geometry admitting a full \emph{analytic} solution. In general fluid dynamics one is forced to numerics; here we get an exact answer that builds intuition.

There are two natural frames, and keeping them straight is essential.

\dfn[fluid-frame]{Fluid Frame (lab frame)}{%
  The frame in which the fluid is at rest at infinity. The boundary condition at infinity is
  \begin{equation}
    \vel \xrightarrow{\;r\to\infty\;} \bvec{0}.
    \label{eq:lec16_bc_fluidframe}
  \end{equation}
  The sphere translates through this frame with velocity $\bvec{v}_0$.
}
\textit{Significance:} This is the ``true'' frame -- far from the sphere the water is genuinely still. There is \emph{no} incoming stream; the only disturbance is what the sphere stirs up locally.

\dfn[sphere-frame]{Sphere Frame}{%
  The frame instantaneously co-moving with the sphere, in which the sphere is at rest and the fluid streams past. The boundary condition at infinity becomes
  \begin{equation}
    \vel \xrightarrow{\;r\to\infty\;} -\bvec{v}_0,
    \label{eq:lec16_bc_sphereframe}
  \end{equation}
  i.e.\ a uniform stream of velocity $-\bvec{v}_0$.
}
\textit{Significance:} This is the ``wind-tunnel'' frame: the aircraft sits still and giant fans blow air past it. It is the more convenient frame to solve in, because the geometry (a fixed sphere in a uniform stream) is static. The asterisk we attach to the name signals a caveat: if the sphere accelerates, this frame is \emph{non-inertial}; it is genuinely inertial only at the instant we freeze the velocity.

The two frames are related by a Galilean boost. If $\phi$ and $\phi^{\text{(fluid)}}$ denote the potentials in the sphere and fluid frames respectively, then since velocities add as $\vel^{\text{(fluid)}}=\vel+\bvec{v}_0$,
\begin{equation}
  \phi^{\text{(fluid)}}(\curcoord,t) = \phi(\curcoord,t) + \bvec{v}_0\cdot\curcoord,
  \label{eq:lec16_galilean}
\end{equation}
because $\grad(\bvec{v}_0\cdot\curcoord)=\bvec{v}_0$ reproduces the extra uniform velocity. We will solve in the sphere frame and then add the $\bvec{v}_0\cdot\curcoord$ piece to descend to the fluid frame.

\begin{figure}[h]
\centering
\begin{tikzpicture}[scale=1.0]
  % ===== Panel A: fluid frame =====
  \begin{scope}[shift={(0,0)}]
    \node[font=\bfseries] at (0,3.2) {Panel A: Fluid Frame};
    \node[font=\small] at (0,2.75) {(fluid at rest at $\infty$)};
    % sphere
    \shade[ball color=myblue!40] (0,0) circle (0.55);
    \node[font=\small] at (0,0) {$R$};
    % velocity of sphere
    \draw[force,line width=1.2pt] (0,0) -- (1.5,0) node[right,font=\small]{$\bvec{v}_0$};
    % dipole streamlines (closed loops pushed ahead/behind)
    \foreach \s in {0.85,1.25,1.7}{
      \draw[myblue,thick] (0,\s) .. controls (1.1*\s,0.4*\s) and (1.1*\s,-0.4*\s) .. (0,-\s);
      \draw[myblue,thick] (0,\s) .. controls (-1.1*\s,0.4*\s) and (-1.1*\s,-0.4*\s) .. (0,-\s);
    }
    \node[font=\small] at (0,-2.6) {closed dipole loops};
  \end{scope}
  % ===== Panel B: sphere frame =====
  \begin{scope}[shift={(6.2,0)}]
    \node[font=\bfseries] at (0,3.2) {Panel B: Sphere Frame};
    \node[font=\small] at (0,2.75) {(uniform stream $-\bvec{v}_0$)};
    \shade[ball color=myblue!40] (0,0) circle (0.55);
    \node[font=\small] at (0,0) {$r=R$};
    % incoming streamlines diverting around the sphere
    \foreach \y in {-1.6,-1.0,1.0,1.6}{
      \draw[vvec] (-2.6,\y) .. controls (-1.0,\y) and (-0.8,{\y*0.55}) .. (0,{\y*0.45})
                  .. controls (0.8,{\y*0.55}) and (1.0,\y) .. (2.6,\y);
    }
    % central streamline hitting stagnation points
    \draw[vvec] (-2.6,0) -- (-0.58,0);
    \draw[vvec] (0.58,0) -- (2.6,0);
    % stagnation points
    \fill[myred] (-0.55,0) circle(2.4pt);
    \fill[myred] (0.55,0) circle(2.4pt);
    \node[myred,font=\scriptsize,below left] at (-0.55,-0.05) {front stag.};
    \node[myred,font=\scriptsize,below right] at (0.55,-0.05) {rear stag.};
    % angle theta
    \draw[->] (0,0) -- (0.39,0.39);
    \draw[->] (0.45,0) arc (0:45:0.45);
    \node[font=\small] at (0.62,0.22) {$\theta$};
  \end{scope}
\end{tikzpicture}
\caption{Two frames for the moving sphere. \textbf{Panel A} (fluid frame): the fluid is still at infinity; the sphere pushes fluid ahead and pulls it in behind, producing closed dipole loops centred on the sphere. \textbf{Panel B} (sphere frame): a uniform stream $-\bvec{v}_0$ diverts around the fixed sphere; the central streamline terminates at the \emph{front} ($\theta=\pi$) and \emph{rear} ($\theta=0$) stagnation points where $\vel=\bvec{0}$.}
\label{fig:lec16_twoframes}
\end{figure}

\subsubsection*{Qualitative Expectations}
Before solving, the lecture pauses to anticipate the answer -- a healthy physicist's habit.

\begin{itemize}
  \item \textbf{Streamlines.} The sphere must push fluid out of its way at the front and let fluid fill in behind it. In the fluid frame this gives the closed loops of Panel~A; in the sphere frame, fluid divides and flows around the obstacle (Panel~B).
  \item \textbf{Stagnation points.} In the sphere frame, the streamline aiming dead-centre at the sphere cannot pass through it; by symmetry it must terminate on the surface. The velocity there is zero -- a \emph{stagnation point}. There are two, fore and aft.
  \item \textbf{Drag?} Everyday experience says yes: a ball dropped in glycerin reaches a terminal velocity, so there must be a retarding \emph{drag force}. We write a question mark next to this expectation -- and we will find, paradoxically, that the ideal-fluid calculation gives \emph{zero drag}.
\end{itemize}

%%─────────────────────────────────────────────────────────────────────────────
\subsection{The Boundary-Value Problem in the Sphere Frame}\label{sec:lec16:bvp}
%%─────────────────────────────────────────────────────────────────────────────

\subsubsection*{Mathematical Setup}
We solve in the sphere frame. The unknown is the velocity potential $\phi$, governed by Laplace's equation \eqref{eq:lec16_laplace}. A unique solution requires boundary conditions at the two boundaries of the domain: infinity and the sphere surface $r=R$.

\paragraph{Far-field boundary condition.}
In the sphere frame the fluid streams uniformly with $\vel\to-\bvec{v}_0$. Integrating $\vel=\grad\phi$ with $\vel=-\bvec{v}_0$ constant:
\begin{equation}
  \phi \xrightarrow{\;r\to\infty\;} -\bvec{v}_0\cdot\curcoord + \text{const}.
  \label{eq:lec16_farfield}
\end{equation}
Choosing the $z$-axis along $\bvec{v}_0$, with $\theta$ the polar angle from $+z$, we have $\bvec{v}_0\cdot\curcoord=v_0\, r\cos\theta$, so
\begin{equation}
  \phi \xrightarrow{\;r\to\infty\;} -v_0\, r\cos\theta.
  \label{eq:lec16_farfield_polar}
\end{equation}
(The additive constant does not change $\vel=\grad\phi$, so we discard it -- a ``gauge'' choice.)

\paragraph{Surface boundary condition (impermeability).}
The sphere is a rigid, impermeable solid: fluid cannot cross its surface, so the \emph{normal} component of velocity must vanish there. The outward normal on a sphere is radial, $\uvec{n}=\uvec{r}$, so the normal velocity is just the radial derivative of $\phi$:
\begin{equation}
  \boxed{v_n = \uvec{n}\cdot\grad\phi = \pdv{\phi}{r}\bigg|_{r=R} = 0.}
  \label{eq:lec16_bc_neumann}
\end{equation}
\textit{Significance:} This is a \textbf{Neumann} boundary condition (a condition on the normal derivative, not the value, of $\phi$). It encodes the single physical fact that the wall is solid. This contrasts with the analogous electrostatics problem, where a conductor imposes a \emph{Dirichlet} condition ($\phi_{\text{elec}}=\text{const}$).

\subsubsection{Electrostatic Analogy}

\subsubsection*{Conceptual Background}
The problem ``$\lapl\phi=0$, with $\grad\phi\to$ const at infinity and a sphere as inner boundary'' is one you have already solved in electromagnetism: a \emph{conducting sphere in a uniform external field $\bvec{E}_0$}. There, the sphere distorts the field and develops an \emph{induced dipole}. The two problems share the same equation and the same far-field structure ($\grad\phi\to$ const); they differ \emph{only} in the surface condition:
\begin{center}
\begin{tabular}{lll}
  \toprule
  & Sphere in flow & Conducting sphere in $\bvec{E}_0$ \\
  \midrule
  Equation & $\lapl\phi=0$ & $\lapl\phi_{\text{el}}=0$ \\
  At $\infty$ & $\grad\phi\to-\bvec{v}_0$ & $\grad\phi_{\text{el}}\to-\bvec{E}_0$ \\
  At $r=R$ & $\partial_r\phi=0$ (Neumann) & $\phi_{\text{el}}=\text{const}$ (Dirichlet) \\
  \bottomrule
\end{tabular}
\end{center}
D'Alembert solved this flow problem in 1750, long before electrostatics existed; the modern reader gets the solution almost for free by analogy. In both cases the response of the sphere is a pure \emph{dipole}.

\subsubsection{Multipole Expansion and Selection of the Dipole}

\subsubsection*{Mathematical Setup}
The general axisymmetric solution of Laplace's equation, written as a sum over spherical harmonics $Y_{\ell m}$, is
\begin{equation}
  \phi(r,\theta,\varphi) = \sum_{\ell,m}\left(a_{\ell m}\,r^{\ell} + \frac{b_{\ell m}}{r^{\ell+1}}\right)Y_{\ell m}(\theta,\varphi).
  \label{eq:lec16_multipole}
\end{equation}
For each multipole $\ell$ the radial dependence is a pure power law: a \emph{growing} solution $r^{\ell}$ (sourced from infinity) and a \emph{decaying} solution $r^{-(\ell+1)}$ (sourced from the origin). The angular dependence is carried entirely by the spherical harmonics.

\paragraph{Why a pure dipole?}
The far-field condition \eqref{eq:lec16_farfield_polar} has the angular dependence $\cos\theta$. But $\cos\theta$ \emph{is} (proportional to) the $\ell=1$, $m=0$ spherical harmonic -- the \textbf{dipole}. Listing the lowest harmonics:
\begin{equation}
  Y_{0}\propto 1\ (\text{monopole}),\qquad Y_{1}\propto\cos\theta\ (\text{dipole}),\qquad Y_2\propto(3\cos^2\theta-1)\ (\text{quadrupole}).
  \label{eq:lec16_harmonics}
\end{equation}
The far field is purely $\ell=1$. Because the surface condition is imposed at a \emph{spherical} boundary $r=R$ -- which respects full rotational symmetry -- it cannot mix different $\ell$ values: a $\cos\theta$ forcing can only excite a $\cos\theta$ response. Hence the \emph{entire} solution lives in the $\ell=1$ sector. We therefore take the trial form (keeping the growing $r^1$ piece fixed by the far field, plus the decaying $r^{-2}$ dipole tail):
\begin{equation}
  \phi(r,\theta) = \cos\theta\left(a\, r + \frac{b}{r^2}\right) = -v_0\cos\theta\left(r + \frac{b}{r^2}\right),
  \label{eq:lec16_trial}
\end{equation}
where matching \eqref{eq:lec16_farfield_polar} fixes the growing coefficient to $a=-v_0$ (so that as $r\to\infty$ only the $-v_0 r\cos\theta$ survives). The single remaining unknown is the dipole strength $b$.

\textit{Significance:} A \emph{vector} boundary datum at infinity (the uniform velocity $\bvec{v}_0$, a directional object) maps onto exactly the dipole harmonic, because a dipole is the unique $\ell=1$ object and $\ell=1$ is the unique multipole carrying a single preferred direction. This is why ``uniform stream past a sphere'' is synonymous with ``pure dipole.''

\subsubsection{Matching the Neumann Condition}

\subsubsection*{Mathematical Setup}
Impose \eqref{eq:lec16_bc_neumann} on the trial \eqref{eq:lec16_trial}. The normal derivative is the radial derivative:
\begin{equation}
  \pdv{\phi}{r} = -v_0\cos\theta\left(1 - \frac{2b}{r^3}\right).
  \label{eq:lec16_dphidr}
\end{equation}
Evaluate at $r=R$ and demand it vanish for \emph{all} $\theta$:
\begin{equation}
  \pdv{\phi}{r}\bigg|_{r=R} = -v_0\cos\theta\left(1 - \frac{2b}{R^3}\right) = 0.
  \label{eq:lec16_neumann_match}
\end{equation}
Since $\cos\theta$ is not identically zero, the bracket must vanish:
\begin{equation}
  1 - \frac{2b}{R^3} = 0 \;\Longrightarrow\; b = \frac{R^3}{2}.
  \label{eq:lec16_solve_b}
\end{equation}
\textit{Significance:} The condition collapses to a single algebraic equation for $b$ precisely because the $\cos\theta$ factors out of every term -- a direct dividend of the multipole selection. There is exactly one dipole strength that makes the sphere impermeable.

\subsubsection{The Solution in the Sphere Frame}

Substituting $b=R^3/2$ into \eqref{eq:lec16_trial}:
\begin{equation}
  \boxed{\phi(r,\theta) = -v_0\cos\theta\left(r + \frac{R^3}{2r^2}\right) \quad\text{(sphere frame)}.}
  \label{eq:lec16_phi_sphereframe}
\end{equation}
\textit{Significance:} This is the exact velocity potential for uniform flow past a sphere. The first term is the undisturbed stream; the second is the dipole correction the sphere adds to enforce impermeability. One checks the two boundary conditions directly: as $r\to\infty$ the dipole term dies and $\phi\to-v_0 r\cos\theta$; at $r=R$, $\partial_r\phi\propto(1-2b/R^3)=0$, so no fluid penetrates the surface.

%%─────────────────────────────────────────────────────────────────────────────
\subsection{Transforming to the Fluid Frame}\label{sec:lec16:fluidframe}
%%─────────────────────────────────────────────────────────────────────────────

\subsubsection*{Mathematical Setup}
To get the potential in the (lab) fluid frame we apply the Galilean rule \eqref{eq:lec16_galilean}, adding $\bvec{v}_0\cdot\curcoord=v_0 r\cos\theta$:
\begin{equation}
  \phi^{\text{(fluid)}} = \phi + v_0 r\cos\theta
  = -v_0\cos\theta\left(r + \frac{R^3}{2r^2}\right) + v_0 r\cos\theta.
  \label{eq:lec16_fluid_add}
\end{equation}
The two uniform-stream pieces $\mp v_0 r\cos\theta$ cancel exactly, leaving only the dipole:
\begin{equation}
  \boxed{\phi^{\text{(fluid)}}(r,\theta) = -\frac{v_0 R^3\cos\theta}{2r^2} \quad\text{(fluid frame)}.}
  \label{eq:lec16_phi_fluidframe}
\end{equation}
\textit{Significance:} In the fluid frame the potential is a \emph{pure dipole} -- nothing else. The growing $r^1$ ``incoming-from-infinity'' multipole has been removed, consistent with the physical fact that there is no flow at infinity in this frame. Of the two kinds of multipole in \eqref{eq:lec16_multipole}, the fluid frame retains only the $1/r^2$ dipole \emph{sourced at the origin} (decaying outward); the growing dipole ``arriving from infinity,'' which was present in the sphere frame as the uniform stream, has cancelled.

\nt{Multipoles come in two species: those \emph{sourced at the origin}, decaying as $r^{-(\ell+1)}$, and those \emph{arriving from infinity}, growing as $r^{\ell}$. A monopole field goes as $1/r^2$ in velocity, a dipole as $1/r^3$. The fluid-frame solution is the cleanest statement of the physics: a moving sphere is, to the surrounding still fluid, exactly a hydrodynamic dipole.}

%%─────────────────────────────────────────────────────────────────────────────
\subsection{Streamlines, Equipotentials, and Stagnation Points}\label{sec:lec16:streamlines}
%%─────────────────────────────────────────────────────────────────────────────

\subsubsection*{Physical Motivation}
With $\phi$ in hand we can verify the qualitative picture of \cref{fig:lec16_twoframes} and locate the stagnation points quantitatively. The velocity components in spherical coordinates follow from $\vel=\grad\phi$:
\begin{equation}
  v_r = \pdv{\phi}{r},\qquad v_\theta = \frac{1}{r}\pdv{\phi}{\theta},\qquad v_\varphi = \frac{1}{r\sin\theta}\pdv{\phi}{\varphi}=0,
  \label{eq:lec16_vel_components}
\end{equation}
the azimuthal component vanishing because the solution is axisymmetric ($\partial_\varphi\phi=0$).

\subsubsection{Stagnation Points}

\subsubsection*{Mathematical Setup}
Work in the sphere frame, \eqref{eq:lec16_phi_sphereframe}. The tangential component is
\begin{equation}
  v_\theta = \frac{1}{r}\pdv{\phi}{\theta} = \frac{1}{r}\,v_0\sin\theta\left(r+\frac{R^3}{2r^2}\right)\propto\sin\theta,
  \label{eq:lec16_vtheta}
\end{equation}
which vanishes at $\theta=0$ and $\theta=\pi$ -- the poles. The radial component is
\begin{equation}
  v_r = \pdv{\phi}{r} = -v_0\cos\theta\left(1-\frac{R^3}{r^3}\right),
  \label{eq:lec16_vr}
\end{equation}
which we \emph{constructed} to vanish on the surface $r=R$ (the bracket is $1-R^3/r^3$). Therefore at the two surface poles $(r=R,\ \theta=0)$ and $(r=R,\ \theta=\pi)$ \emph{both} components vanish (and $v_\varphi=0$ always), so
\begin{equation}
  \boxed{\vel = \bvec{0}\quad\text{at the front and rear poles of the sphere.}}
  \label{eq:lec16_stagnation}
\end{equation}
\textit{Significance:} These are the two \emph{stagnation points} anticipated qualitatively. The rear pole (in the direction of motion) sits at $\theta=0$, the front pole at $\theta=\pi$. By Bernoulli \eqref{eq:lec16_pressure}, $v=0$ means \emph{maximum} pressure: the fluid presses hardest exactly fore and aft. The perfect fore--aft symmetry of these high-pressure zones is the seed of D'Alembert's paradox.

\subsubsection{Equipotential Surfaces}

In the fluid frame, the dipole equipotentials $\phi^{\text{(fluid)}}=\text{const}$ satisfy, from \eqref{eq:lec16_phi_fluidframe},
\begin{equation}
  \frac{\cos\theta}{r^2}=\text{const}\;\Longrightarrow\; r^2 = r_{\max}^2\,\cos\theta,
  \label{eq:lec16_equipotential}
\end{equation}
where $r_{\max}$ is the radius reached at $\theta=0$. Each such surface passes through the origin (at $\theta=\pi/2$, $r\to 0$) and is a closed lobe; different constants give nested, scaled copies of the same lobe, with a mirror family for $\cos\theta<0$. The full surface is the body of revolution obtained by rotating the lobe about the $z$-axis (the motion axis), giving a closed dimpled surface that pinches to a singular point at the origin. The streamlines (blue in \cref{fig:lec16_twoframes}) are everywhere perpendicular to these equipotentials, since $\vel=\grad\phi$.

\nt{In lecture the professor first mis-drew the equipotential surface as a torus; the correct surface is the figure of revolution about the \emph{motion axis}, a closed lobe pinched at the origin, not a doughnut. The streamline picture and all formulae were correct.}

%%─────────────────────────────────────────────────────────────────────────────
\subsection{Generalization to Arbitrary Motion}\label{sec:lec16:general}
%%─────────────────────────────────────────────────────────────────────────────

\subsubsection*{Conceptual Background}
So far the sphere moved at \emph{constant} velocity. We now allow the sphere centre to follow an arbitrary trajectory $\bvec{r}_s(t)$ with instantaneous velocity $\dot{\bvec{r}}_s(t)$. The key structural observation is that \emph{time does not appear in Laplace's equation} \eqref{eq:lec16_laplace} -- it enters only through the boundary conditions. So at each instant the spatial problem is the \emph{same} dipole problem, with two updates: (i) re-centre the coordinates on the moving sphere, $\curcoord\to\curcoord-\bvec{r}_s(t)$; (ii) replace the constant $\bvec{v}_0$ by the instantaneous velocity $\dot{\bvec{r}}_s(t)$.

\subsubsection{Vectorial Rewriting}

\subsubsection*{Mathematical Setup}
First rewrite the fluid-frame dipole \eqref{eq:lec16_phi_fluidframe} covariantly. Note $v_0\cos\theta = \bvec{v}_0\cdot\uvec{r}$, hence
\begin{equation}
  \frac{v_0\cos\theta}{r^2} = \frac{\bvec{v}_0\cdot\uvec{r}}{r^2} = \frac{\bvec{v}_0\cdot\curcoord}{r^3},
  \label{eq:lec16_vector_rewrite}
\end{equation}
using $\uvec{r}=\curcoord/r$. Thus the constant-velocity fluid-frame potential is
\begin{equation}
  \phi^{\text{(fluid)}} = -\frac{R^3}{2}\,\frac{\bvec{v}_0\cdot\curcoord}{r^3}.
  \label{eq:lec16_phi_vector}
\end{equation}

\subsubsection{The Accelerating Sphere}

Now promote to arbitrary motion. Measuring position relative to the instantaneous sphere centre, $\curcoord\to\curcoord-\bvec{r}_s(t)$, and the velocity to the instantaneous value $\dot{\bvec{r}}_s(t)$:
\begin{equation}
  \boxed{\phi(\curcoord,t) = -\frac{R^3}{2}\,\frac{\dot{\bvec{r}}_s(t)\cdot\big(\curcoord-\bvec{r}_s(t)\big)}{\,\lvert\curcoord-\bvec{r}_s(t)\rvert^{3}}.}
  \label{eq:lec16_phi_general}
\end{equation}
\textit{Significance:} This is the exact potential for a sphere in \emph{arbitrary} translational motion through ideal fluid. The crucial point: the surface boundary condition $v_n=0$ depends only on the instantaneous \emph{velocity} of the sphere, never on its acceleration -- a solid wall is impermeable regardless of how fast it accelerates. Acceleration enters the physics only later, through the \emph{time derivative} $\partial_t\phi$ in the pressure \eqref{eq:lec16_pressure}, never through the geometry of $\phi$ itself. This is exactly why the static-looking dipole suffices even for accelerated motion.

%%─────────────────────────────────────────────────────────────────────────────
\subsection{The Force on the Sphere}\label{sec:lec16:force}
%%─────────────────────────────────────────────────────────────────────────────

\subsubsection*{Conceptual Background}
How does the fluid exert a force on the sphere? Exactly as in elasticity: neighbouring fluid elements exert short-range \emph{contact} stresses on one another, described by the \emph{stress tensor} $\stress$. In an ideal fluid the stress has no deviatoric (shear) part -- there is no viscosity -- so it is purely isotropic:
\begin{equation}
  \sigma_{ij} = -p\,\delta_{ij}.
  \label{eq:lec16_stress_ideal}
\end{equation}
The coefficient of the isotropic stress, with the conventional minus sign, is the \emph{pressure}. The total force the fluid exerts on the sphere is the surface integral of this stress over the sphere's surface. With the area element $d\bvec{S}=\uvec{n}\,dS$ pointing \emph{outward} (out of the fluid, into the sphere is $-\uvec{n}$), the force on the sphere is
\begin{equation}
  \boxed{\bvec{F} = -\oint_{r=R} p\,\uvec{n}\,dS = -\oint_{r=R} p(\theta,\varphi)\,\uvec{r}\,R^2\,d\Omega.}
  \label{eq:lec16_force_def}
\end{equation}
\textit{Significance:} Drag and lift are nothing more than the imbalance of pressure pushing on different parts of the surface. If the pressure pattern is fore--aft symmetric, the net push is zero. The whole question of drag thus reduces to the \emph{symmetry} of the pressure field.

\subsubsection{Assembling the Pressure}

\subsubsection*{Mathematical Setup}
From the unsteady Bernoulli relation \eqref{eq:lec16_pressure} (incompressible, gravity dropped),
\begin{equation}
  p = \rho f(t) - \rho\,\pdv{\phi}{t} - \tfrac{1}{2}\rho v^2.
  \label{eq:lec16_p_three}
\end{equation}
There are three contributions. We examine which survive the surface integral \eqref{eq:lec16_force_def}. The decisive tool is the \textbf{parity symmetry} $\curcoord\to-\curcoord$ about the sphere centre: it maps each surface point to its antipode. A term in $p$ contributes \emph{zero net force} if it produces \emph{equal} pressure at antipodal points, because the outward normals $\uvec{r}$ and $-\uvec{r}$ then cancel in the integral.

\paragraph{(1) The $\rho f(t)$ term.}
Spatially uniform: equal at all points, hence equal at antipodes. It pushes inward equally from all directions and integrates to zero force,
\begin{equation}
  -\oint \rho f(t)\,\uvec{r}\,R^2\,d\Omega = -\rho f(t)R^2\oint\uvec{r}\,d\Omega = \bvec{0},
  \label{eq:lec16_ft_zero}
\end{equation}
since $\oint\uvec{r}\,d\Omega=\bvec{0}$ (the average of the unit normal over a sphere vanishes).

\paragraph{(2) The $-\tfrac{1}{2}\rho v^2$ term.}
The velocity field is dipolar, like the field of an electric dipole, so it has the property $\bvec{v}(-\curcoord)=-\bvec{v}(\curcoord)$ (odd under parity). Therefore its square is \emph{even}:
\begin{equation}
  v^2(-\curcoord) = v^2(\curcoord).
  \label{eq:lec16_vsq_even}
\end{equation}
Equal pressure at antipodal points $\Rightarrow$ zero net force. \emph{The Bernoulli suction does not produce drag.}

\paragraph{(3) The unsteady term $-\rho\,\partial_t\phi$.}
This is the only term that can survive. Differentiate \eqref{eq:lec16_phi_general}. Two pieces depend on time: the explicit velocity $\dot{\bvec{r}}_s$ and the position $\bvec{r}_s$. Writing $\bvec{r}'=\curcoord-\bvec{r}_s$,
\begin{equation}
  \pdv{\phi}{t} = -\frac{R^3}{2}\left[\frac{\ddot{\bvec{r}}_s\cdot\bvec{r}'}{r'^3}
  + \dot{\bvec{r}}_s\cdot\pdv{t}\!\left(\frac{\bvec{r}'}{r'^3}\right)\right].
  \label{eq:lec16_dtphi}
\end{equation}
The first piece is proportional to the \emph{acceleration} $\ddot{\bvec{r}}_s$; under parity $\bvec{r}'\to-\bvec{r}'$ it is \emph{odd}, so it gives equal-and-opposite -- actually \emph{unequal} -- pressure at antipodes and can yield a nonzero force. The second piece, $\dot{\bvec{r}}_s\cdot\partial_t(\bvec{r}'/r'^3)$, is built from velocities only; since $\partial_t(\bvec{r}'/r'^3)=-(\dot{\bvec{r}}_s\cdot\grad)(\bvec{r}'/r'^3)$ it is quadratic in velocity and even under parity, contributing zero force by the same argument as term~(2).

\textit{Significance:} Of the three pressure contributions, only the piece proportional to the sphere's \emph{acceleration} can exert a net force. In particular, \textbf{at constant velocity ($\ddot{\bvec{r}}_s=\bvec{0}$) every term integrates to zero -- the net force vanishes.}

\thm[no-drag]{D'Alembert's Paradox (No Steady Drag)}{%
  An ideal (inviscid, incompressible, irrotational) fluid exerts \emph{zero drag} on a sphere moving at constant velocity:
  \begin{equation}
    \boxed{\bvec{F}_{\text{drag}} = \bvec{0}\quad\text{for}\quad\dot{\bvec{r}}_s = \text{const}.}
    \label{eq:lec16_dalembert}
  \end{equation}
}
\textit{Significance:} Drag, by definition, is a force opposing steady motion. The ideal-fluid model predicts \emph{none} -- in flat contradiction with the everyday observation that objects slow down in fluids. This is the famous paradox D'Alembert reached in 1750. Its root is the perfect fore--aft symmetry of the pressure field (Panel~A): the high pressure at the front stagnation point is exactly matched by equal high pressure at the rear, so the pushes cancel.

%%─────────────────────────────────────────────────────────────────────────────
\subsection{Added Mass (Effective Mass)}\label{sec:lec16:addedmass}
%%─────────────────────────────────────────────────────────────────────────────

\subsubsection*{Physical Motivation}
The surviving acceleration term does produce a force. Let us compute it. Physically, an accelerating sphere must \emph{drag the surrounding fluid along with it}; setting that fluid into motion costs an extra force beyond what is needed to accelerate the bare sphere. The sphere therefore behaves as if it were heavier -- it has an \emph{added mass}.

\subsubsection{A Tensor Integral over the Sphere}

\subsubsection*{Mathematical Setup}
The force \eqref{eq:lec16_force_def} from the acceleration term requires the angular integral of $\hat r_i \hat r_j$ over the sphere. We compute it once and for all.

\lem[rr-integral]{Angular Average of $\hat r_i\hat r_j$}{%
  Over a sphere of radius $R$,
  \begin{equation}
    \oint \hat r_i\,\hat r_j\,d\Omega = \frac{4\pi}{3}\,\delta_{ij}.
    \label{eq:lec16_rr_int}
  \end{equation}
}
\begin{proof}
By rotational invariance the integral can only depend on the one available isotropic tensor, $\delta_{ij}$: write $\oint \hat r_i\hat r_j\,d\Omega = C\,\delta_{ij}$. Take the trace ($i=j$, sum): the left side becomes $\oint \hat r_i\hat r_i\,d\Omega = \oint 1\,d\Omega = 4\pi$ (since $\hat r_i\hat r_i=|\uvec{r}|^2=1$), while the right side is $C\,\delta_{ii}=3C$. Hence $3C=4\pi$, giving $C=4\pi/3$.
\end{proof}
\textit{Significance:} This is the workhorse identity: the angular average of two unit-normal components is isotropic, with coefficient $4\pi/3$. It is exactly what converts the dipolar pressure pattern into a clean force along the acceleration direction.

\subsubsection{Evaluating the Force}

\subsubsection*{Mathematical Setup}
Keep only the surviving acceleration piece of the pressure. From \eqref{eq:lec16_p_three} and the first term of \eqref{eq:lec16_dtphi}, on the surface $r'=R$,
\begin{equation}
  p\big|_{\text{accel}} = -\rho\,\pdv{\phi}{t}\bigg|_{\text{accel}} = +\frac{\rho R^3}{2}\,\frac{\ddot{\bvec{r}}_s\cdot\uvec{r}}{R^2}
  = \frac{\rho R}{2}\,\ddot r_{s,j}\,\hat r_j.
  \label{eq:lec16_p_accel}
\end{equation}
Insert into the force \eqref{eq:lec16_force_def}, with $dS=R^2 d\Omega$ and $\uvec{n}=\uvec{r}$:
\begin{equation}
  F_i = -\oint p\,\hat r_i\,R^2\,d\Omega
  = -\frac{\rho R}{2}\,\ddot r_{s,j}\,R^2\oint \hat r_i\,\hat r_j\,d\Omega.
  \label{eq:lec16_F_setup}
\end{equation}
Apply \eqref{eq:lec16_rr_int}:
\begin{equation}
  F_i = -\frac{\rho R^3}{2}\,\ddot r_{s,j}\cdot\frac{4\pi}{3}\,\delta_{ij}
  = -\frac{4\pi}{3}\cdot\frac{\rho R^3}{2}\,\ddot r_{s,i}.
  \label{eq:lec16_F_eval}
\end{equation}
Recognise $V_{\text{sph}}=\tfrac{4}{3}\pi R^3$ as the volume of the sphere:
\begin{equation}
  \boxed{\bvec{F} = -\tfrac{1}{2}\,\rho\,V_{\text{sph}}\,\ddot{\bvec{r}}_s.}
  \label{eq:lec16_force_addedmass}
\end{equation}
\textit{Significance:} The fluid pushes \emph{against} the acceleration of the sphere (note the minus sign), with magnitude $\tfrac{1}{2}\rho V_{\text{sph}}$ times the acceleration. It behaves exactly like an extra inertia. There is \emph{no} velocity-dependent (drag) term -- consistent with \cref{th:no-drag}.

\subsubsection{The Effective Mass}

\subsubsection*{Conceptual Background}
Put this force into Newton's second law for the sphere of true mass $m$ under an external force $\bvec{F}_{\text{ext}}$:
\begin{equation}
  m\,\ddot{\bvec{r}}_s = \bvec{F}_{\text{fluid}} + \bvec{F}_{\text{ext}}
  = -\tfrac{1}{2}\rho V_{\text{sph}}\,\ddot{\bvec{r}}_s + \bvec{F}_{\text{ext}}.
  \label{eq:lec16_newton}
\end{equation}
Move the fluid force to the left:
\begin{equation}
  \left(m + \tfrac{1}{2}\rho V_{\text{sph}}\right)\ddot{\bvec{r}}_s = \bvec{F}_{\text{ext}}.
  \label{eq:lec16_newton_eff}
\end{equation}
This is Newton's law with a modified inertia:

\dfn[eff-mass]{Effective (Added) Mass of a Sphere}{%
  A sphere accelerating through ideal fluid responds to external force as if it had mass
  \begin{equation}
    \boxed{m_{\text{eff}} = m + \tfrac{1}{2}\rho V_{\text{sph}} = m + \tfrac{1}{2}m_{\text{displaced}}.}
    \label{eq:lec16_meff}
  \end{equation}
}
\textit{Significance:} The sphere ``feels heavier'' by exactly \emph{half} the mass of the fluid it displaces. The factor $\tfrac{1}{2}$ is specific to the sphere; it comes from the dipole field the sphere generates. The physical origin: to accelerate, the sphere must also accelerate the cloud of fluid moving with it, and that co-moving fluid stores extra kinetic energy.

\nt{Compare with buoyancy (Lec~11): there, the displaced fluid mass appears in a \emph{static} force (Archimedes). Here the displaced fluid mass appears in a \emph{dynamic}, inertial role -- and is multiplied by $\tfrac{1}{2}$. Same $\rho V$, completely different physics: buoyancy is gravitational, added mass is inertial.}

\subsubsection{Generalization: The Effective-Mass Tensor}

\subsubsection*{Conceptual Background}
For a non-spherical body the co-moving fluid -- and hence the added inertia -- depends on the \emph{direction} of acceleration: a flat plate accelerated broadside drags far more fluid than the same plate accelerated edge-on. The added mass is therefore a \textbf{symmetric rank-2 tensor} $M_{ij}$:
\begin{equation}
  F_i = -M_{ij}\,\ddot r_{s,j},\qquad M_{ij}=M_{ji}.
  \label{eq:lec16_Mtensor}
\end{equation}
Its symmetry is most transparent from the kinetic energy of the entrained fluid, which is a quadratic form in the body velocity:
\begin{equation}
  \boxed{T_{\text{fluid}} = \tfrac{1}{2}\,M_{ij}\,v_i\,v_j.}
  \label{eq:lec16_Tfluid}
\end{equation}
\textit{Significance:} Instead of a scalar mass multiplying $v^2$, a general body carries a \emph{matrix} of masses contracting $v_iv_j$. For the sphere, isotropy forces $M_{ij}=\tfrac{1}{2}\rho V_{\text{sph}}\,\delta_{ij}$, recovering the scalar result \eqref{eq:lec16_meff}. The persistence of $\bvec{F}_{\text{drag}}=\bvec{0}$, however, is general: any body with reflection symmetry exhibits zero ideal-fluid drag.

\nt{The general added mass is of order $\rho\times(\text{body volume})$ but the precise coefficient depends on the flow field the body's shape generates -- for the sphere it happens to be exactly $\tfrac{1}{2}$.}

%%─────────────────────────────────────────────────────────────────────────────
\subsection{Resolution of D'Alembert's Paradox}\label{sec:lec16:resolution}
%%─────────────────────────────────────────────────────────────────────────────

\subsubsection*{Conceptual Background}
A paradox in physics is a productive crisis: two strong arguments collide, and resolving the collision teaches something deep. Here, Euler's ideal-flow equations \emph{rigorously} predict zero drag (\cref{th:no-drag}), while everyday experience -- a ball thrown into water decelerating to terminal velocity -- \emph{demands} drag. One of the two must be giving way, and we know which.

The culprit is the \textbf{ideal-fluid approximation itself}. We assumed the stress was purely isotropic, $\sigma_{ij}=-p\,\delta_{ij}$, with \emph{no} shear (viscous) part. Real fluids have \emph{viscosity}: tangential stresses that act at the body surface, enforce a no-slip boundary condition, and -- crucially -- break the fore--aft symmetry of the flow by generating a wake behind the body. Once the viscous shear stress is restored to the stress tensor (the Navier--Stokes equations), the pressure symmetry is destroyed and a genuine drag force appears.

\clm[resolution]{Viscosity Resolves the Paradox}{%
  Drag is absent in ideal flow precisely because the deviatoric (shear) stress was discarded. Restoring viscosity -- moving from Euler to Navier--Stokes -- breaks the fore--aft symmetry, produces a wake, and yields a nonzero drag force consistent with experiment.
}
\textit{Significance:} D'Alembert (1750) could not have spotted this: the concept of viscosity as a momentum-diffusing shear stress was not formulated until Navier and Stokes nearly a century later. The paradox is thus a precise diagnostic -- it isolates \emph{exactly} the term (viscous shear) whose omission renders the model unphysical, and motivates the viscous-flow theory we turn to next.

%%─────────────────────────────────────────────────────────────────────────────
\subsection*{Summary}
%%─────────────────────────────────────────────────────────────────────────────

\begin{itemize}
  \item \textbf{Setup.} A sphere moving through still fluid; solve $\lapl\phi=0$ with $\partial_r\phi|_{R}=0$ (Neumann, impermeable wall) and $\phi\to-v_0 r\cos\theta$ at infinity (uniform stream in sphere frame).
  \item \textbf{Multipole selection.} The $\cos\theta$ far field is the $\ell=1$ harmonic; the spherical boundary cannot mix multipoles, so the solution is a \emph{pure dipole}. Dipole strength fixed by Neumann: $b=R^3/2$.
  \item \textbf{Solutions.} Sphere frame $\phi=-v_0\cos\theta\,(r+R^3/2r^2)$; fluid frame $\phi=-v_0R^3\cos\theta/2r^2$ (pure dipole). Related by Galilean boost $\phi^{\text{fluid}}=\phi+\bvec{v}_0\cdot\curcoord$.
  \item \textbf{Stagnation points.} $\vel=\bvec{0}$ at front and rear poles; pressure maximal there.
  \item \textbf{Arbitrary motion.} $\phi=-\tfrac{R^3}{2}\,\dot{\bvec{r}}_s\!\cdot\!(\curcoord-\bvec{r}_s)/|\curcoord-\bvec{r}_s|^3$; boundary condition needs only the instantaneous velocity, never acceleration.
  \item \textbf{Force.} From $p=\rho f(t)-\rho\partial_t\phi-\tfrac12\rho v^2$: the uniform and $v^2$ terms cancel by parity; only the acceleration term survives. \textbf{Zero drag} at constant velocity (D'Alembert's paradox).
  \item \textbf{Added mass.} $\bvec{F}=-\tfrac12\rho V_{\text{sph}}\,\ddot{\bvec{r}}_s$, so $m_{\text{eff}}=m+\tfrac12\rho V_{\text{sph}}$ (half the displaced fluid mass). Generalizes to a symmetric tensor $M_{ij}$ with $T_{\text{fluid}}=\tfrac12 M_{ij}v_iv_j$.
  \item \textbf{Resolution.} The paradox is cured by viscosity: restoring shear stress (Navier--Stokes) breaks fore--aft symmetry and produces real drag.
\end{itemize}

\subsection*{Further Reading}
\begin{itemize}
  \item L.\,D. Landau and E.\,M. Lifshitz, \textit{Fluid Mechanics} (Vol.\ 6): \S10--11 (flow past a sphere, D'Alembert's paradox, induced/added mass).
  \item D.\,J. Acheson, \textit{Elementary Fluid Dynamics}: Ch.\ 4--5 (potential flow past a sphere, added mass, D'Alembert).
  \item J.\,D. Jackson, \textit{Classical Electrodynamics}: \S2.5 (conducting sphere in a uniform field -- the mathematically identical problem).
\end{itemize}

\end{document}
